\subsubsection{LLMListPage}
\textbf{Modulo}
\begin{itemize}
    \item LLMListModule, modulo funzionale dell'applicazione che offre i servizi di collegamento legati alla creazione, visualizzazione e cancellazione degli LLM. 
\end{itemize}
\textbf{Elementi chiave}
\begin{itemize}
    \item lista scorrevole degli LLM salvati, realizzata tramite sottocomponente \textbf{LLMListView};
    \item popup di conferma di eliminazione di un LLM \textbf{ConfirmComponent};
    \item messaggio di errore visualizzato in seguito ad'un errore nell'ottenimento o nel salvataggio dei dati, realizzato tramite sottocomponente \textbf{MessageComponent}.
\end{itemize}
\textbf{Servizio} \\
\begin{itemize}
    \item LLMListService, servizio che permette l'esecuzione di operazioni CRUD sulla lista dei LLM salvati.
\end{itemize}

\subsubsection{LLMListView}
\textbf{Componente padre}
\begin{itemize}
    \item LLMListPage.
\end{itemize}
\textbf{Elementi chiave}
\begin{itemize}
    \item \textbf{LLMList}: lista scorrevole degli LLM salvati nel sistema, per ogni LLM salvato esiste un suo \textbf{DatasetElement}.
    \item \textbf{StatusMessage}: visualizzato un messaggio che indica all'utente nel caso non ci siano LLM salvati nel sistema;
\end{itemize}
\textbf{Input}
\begin{itemize}
    \item LLM: lista contenente tutti i LLM salvati nel sistema.
\end{itemize}
\textbf{Output}
\begin{itemize}
    \item view LLM: evento emesso quando riceve da \textbf{LLMElement} la richiesta di visuazzazione di un LLM;
    \item add LLM: evento emesso quando riceve da \textbf{LLMElement} la richiesta di aggiunta di un LLM;
    \item delete LLM: evento emesso quando riceve da \textbf{LLMElement} la richiesta di eliminazione di un LLM;
\end{itemize}

\subsubsection{LLMElement}
\textbf{Componente padre}
\begin{itemize}
    \item LLMListView.
\end{itemize}
\textbf{Elementi chiave}
\begin{itemize}
    \item \textbf{viewButton}: pulsante che porta nella pagina di \textbf{LLMPage} per visualizzare in dettaglio del LLM Selezionato.
    \item \textbf{AddButton}: pulsante per aggiungere un nuovo LLM; 
    \item \textbf{DeleteButton}: permette di cancellare il LLM corrispondente.
\end{itemize}
\textbf{Input}
\begin{itemize}
    \item LLM: LLM associato al componente.
\end{itemize}
\textbf{Output}
\begin{itemize}
    \item view: evento emesso quando viene premuto il pulsante per richiedere visualizzazione dell'LLM.
    \item add: evento emesso quando viene premuto il pulsante per richiedere aggiunta di un nuovo LLM;
    \item delete: evento emesso quando viene premuto il pulsante per richiedere eliminazione dell'LLM;
\end{itemize}


\subsubsection{LLMPage}
\textbf{Modulo}
\begin{itemize}
    \item LLMModule, modulo funzionale dell'applicazione che raccoglie e organizza tutti i servizi e risorse legati alla creazione, visualizzazione e modifica di LLM. 
\end{itemize}
\textbf{Elementi chiave}
\begin{itemize}
    \item visualizza le informazioni dettagliate riguardo al LLM selezionato, realizzato tramite sottocomponente \textbf{LLMView} 
    \item visualizza i campi dati da compilare/modificare in caso di una aggiunta LLM, realizzato tramite sottocomponente \textbf{LLMForm}  
    \item messaggio di errore visualizzato in seguito ad'un errore nell'ottenimento o nel salvataggio dei dati, realizzato tramite sottocomponente \textbf{MessageComponent}.
\end{itemize}
\begin{itemize}
    \item LLMService, servizio che permette l'esecuzione di operazioni CRUD sul LLM.
\end{itemize}

\subsubsection{LLMView}
\textbf{Componente padre}
\begin{itemize}
    \item LLMPage.
\end{itemize}
\textbf{Elementi chiave}
\begin{itemize}
    \item \textbf{NameText}: sezione testuale contenente il nome del LLM; 
    \item \textbf{LastUpdateText}: sezione testuale contenente la data dell'ultima modifica del LLM;
    \item \textbf{URLText}: sezione testuale contenente l’URL usato per l’esecuzione delle chiamate HTTP all’API dell’LLM;
    \item \textbf{HeaderText}: sezione testuale contenente le coppie chiave-valore per la costruzione dell’header delle chiamate;
    \item \textbf{BodyText}: sezione testuale contenente le coppie chiave-valore per la costruzione del body delle chiamate;
    \item \textbf{QuestionText}: sezione testuale contenente la chiave da utilizzare per specificare la domanda da porre all’LLM nel body delle richieste HTTP;
    \item \textbf{AnswerText}: sezione testuale contenente la chiave da utilizzare per estrarre la risposta data dall’LLM contenuta nelle risposte HTTP;
    \item \textbf{ModifyButton}: pulsate per abilitare la modifica dell'LLM selezionato.
\end{itemize}
\textbf{input}
\begin{itemize}
    \item selected LLM: viene fornita le informazioni sul LLM selezionata.
\end{itemize}
\textbf{Output}
\begin{itemize}
    \item rename LLM: evento emesso quando viene premuto il pulsante di modificazione;
\end{itemize}

\subsubsection{LLMForm}
\textbf{Componente padre}
\begin{itemize}
    \item LLMPage.
\end{itemize}
\textbf{Elementi chiave}
\begin{itemize}
    \item \textbf{NewNameLLM}: campo di input per modificare il nome dell'LLM;
    \item \textbf{NewURL}: campo di input per modificare L'URL del LLM;
    \item \textbf{NewKeyQuestion}: campo di input per modificare chiave associata alla domanda da porre all'LLM.
    \item \textbf{NewKeyAnswer}: campo di input per modificare chiave associata alla risposta da porre all'LLM.
    \item \textbf{NewKeyHeader} due campi di input per modificare ogni coppia chiave-valore associata alla creazione dell’header delle richieste HTTP verso l’LLM.
    \item \textbf{AddHeaderButton}: pulsante per l'inserimento di una oppia chiave-valore associata alla creazione dell’header delle richieste HTTP verso l’LLM.
    \item \textbf{NewKeyBody} due campi di input per modificare ogni coppia chiave-valore associata alla creazione del body delle richieste HTTP verso l’LLM
    \item \textbf{AddBodyButton}: pulsante per l'inserimento di una oppia chiave-valore associata alla creazione del body delle richieste HTTP verso l’LLM; 
    \item \textbf{SaveButton}: pulsante per salvare le modifiche;
    \item \textbf{AbortButton}: pulsante per annullare la modifica;
    \item \textbf{ValidationFeedback}: messaggi di errore per input vuoti o duplicati;
\end{itemize}
\textbf{input}
\begin{itemize}
    \item selecte: viene fornita le informazioni sul LLM selezionata.
\end{itemize}
\textbf{Output}
\begin{itemize}
    \item confirm: evento emesso quando viene premuto il pulsante di salvataggio;
\end{itemize}